\subsection*{Week of 2026-11-16: polynomials and their roots}

Second rail: recurrences and induction.

\begin{center}\small
\begin{tabular}{@{}c >{\raggedright\arraybackslash}p{3.1cm} >{\raggedright\arraybackslash}p{4.8cm} >{\raggedright\arraybackslash}p{4.8cm}@{}}
\toprule
Rung & Problem & Object and ask & Move \\
\midrule
1 & \putnam{2019 A1} & A symmetric expression in three integers; which values it takes. & Factor the symmetric form. \\
2 & \putnam{2023 A2} & An even-degree polynomial with prescribed values; the remaining solutions. & Build the polynomial from its known values. \\
3 & \seminar{Roots \#6} and \seminar{Roots \#7} & Two Putnam problems from 1968 and 1981 on real roots. & Rolle and sign changes. \\
4 & \putnam{2016 A1} & Derivatives of an integer polynomial; the smallest order that is always divisible by something. & Polynomials meet congruences, as in week 1. \\
5 & \putnam{2022 B1} & A power series built from an integer polynomial; its coefficients are nonzero. & Work modulo 2, which is the week 1 move inside a polynomial. \\
6 & \putnam{2017 A2} & A rational recurrence whose terms turn out to be integer polynomials. & Tagged recurrence; the polynomial identity carries it. Cite the rung where you first wrote a polynomial by its roots or values rather than its coefficients. \\
\midrule
opt & \putnam{2024 A2} & Iterates of a polynomial and divisibility; which polynomials. & \\
opt & \putnam{2019 B5} & A polynomial interpolating a famous recurrence; extrapolate. & \\
\bottomrule
\end{tabular}
\end{center}

Seminar: \seminar{Roots \#1, \#8, \#11a}; \seminar{PS 4 \#1, \#3, \#5, \#9}.
